\section{Broader Impact and Ethics}
\label{sec:ethics}

\paragraph{No human-subjects data.} All fixtures derive from the CDISC
SDTM/ADaM Pilot Project, a public reference dataset distributed by CDISC for
educational use. The dataset contains no real patient records, no protected
health information, and no proprietary trial data. No IRB review is required.

\paragraph{Intended use.} ClinProg-Bench is intended to evaluate code-LLMs
and agentic systems on clinical-programming tasks under reproducible
conditions. It is \emph{not} a certification of fitness for regulatory
submission: a model scoring well on this benchmark is not thereby cleared
for production use in a phase-III trial. Practitioners should continue to
follow ICH E6 (R3), 21 CFR Part 11, and their organisation's validation
SOPs.

\paragraph{Misuse considerations.} The benchmark could in principle be
abused to over-claim AI capability in regulatory contexts. We mitigate this
with: (a) explicit benchmark-vs-deployment scoping in this section,
(b) a deterministic, fully-public evaluation harness that resists cherry-picked
reporting, and (c) a leakage audit and dual-reviewer audit committed to the
repository.

\paragraph{Author qualifications and review.} The dataset was authored by
the listed author with the CDISC Pilot Project as substrate; a three-member
expert panel (clinical statistical programmers in the listed acknowledgements)
audits a stratified subset under a Cohen's $\kappa \geq 0.8$ acceptance
criterion (\S\ref{sec:dataset}).

\paragraph{Carbon footprint.} v0.1 LLM baseline runs (\S\ref{sec:baselines})
total under 5\,M tokens and complete in well under one GPU-hour-equivalent;
the benchmark itself adds negligible compute cost. Future evaluations will
report token counts and wallclock budgets in the leaderboard.

\paragraph{Maintenance.} The author commits to maintaining the benchmark
through the next two NeurIPS review cycles. Release cadence, deprecation
policy, and contact channels are documented in the repository
\texttt{README.md}; the Zenodo DOI provides long-term archival.
